\chapter{Relación del Fenómeno con el Ser Humano}
\label{sec.cap3:ser_humano}

Las consecuencias de un incendio industrial sobre las personas abarcan desde traumatismos y quemaduras térmicas cutáneas
hasta la asfixia celular derivada de los subproductos de la pirólisis. En la práctica pericial y forense se constata que
la inmensa mayoría de las víctimas mortales en incendios estructurales no fallece por la acción directa del fuego o de
las llamas, sino por la inhalación de atmósferas contaminadas con gases tóxicos y partículas en suspensión.

\section{Toxicidad del Monóxido de Carbono y el Cianuro}

Bajo condiciones reales de incendio en recintos cerrados o semi-confinados, la combustión se desarrolla de manera
incompleta debido a la restricción progresiva en el aporte de oxígeno. Esta deficiencia, combinada con la degradación
térmica de los materiales, se liberan concentraciones letales de gases asfixiantes, entre los cuales destacan el
monóxido de carbono y el cianuro de hidrógeno:

\begin{itemize}
    \item Monóxido de Carbono (CO): Es un gas inodoro, incoloro e insípido, lo que impide su percepción sensorial directa
          por los ocupantes. Al ingresar por el tracto respiratorio, atraviesa los alvéolos pulmonares y se enlaza con el
          hierro de la hemoglobina sanguínea con una afinidad aproximadamente $250$ veces superior a la del oxígeno
          molecular. Este fenómeno desplaza al oxígeno y bloquea su transporte. La disminucion progresiva del oxigeno
          produce alteraciones psicomotrices, perdida del conocimiento, coma y colapsos cardiovasculares.
    \item Cianuro de Hidrógeno (HCN): Se desprende durante la descomposición pirolítica de materiales que contienen
          átomos de nitrógeno en su estructura molecular (lanas, poliuretano, nylon y espumas de aislación). El gas
          cianuro es rápidamente absorbido y actúa a nivel celular como un potente inhibidor enzimático, volviendo
          incapaces a las celulas de usar el oxigeno en sangre, produciendo muerte celular.
\end{itemize}

La inhalación combinada de monóxido de carbono y cianuro ejerce un efecto toxicológico sinérgico y multiplicativo. La
primera manifestación clínica consiste en la pérdida de la agudeza visual, confusión mental y desorientación espacial
aguda. Este cuadro anula la capacidad del individuo para orientarse o accionar mecanismos de escape, provocando debilidad
muscular súbita en miembros inferiores que postra a la víctima en el plano del suelo antes de poder alcanzar una salida de
emergencia.

\section{Dinámica del Humo, Pérdida de Visibilidad y Comportamiento Humano}

Además de la asfixia química inducida por monóxido de carbono y cianuro de hidrógeno, el humo de combustión ejerce
mecanismos físicos y psicológicos que degradan la capacidad de supervivencia de los ocupantes:

\begin{itemize}
    \item Gases Irritantes y Corrosivos: La descomposición pirolítica de diversos materiales genera compuestos altamente
          nocivos para las mucosas oculares y el tracto respiratorio:
          \begin{itemize}
              \item Cloruro de hidrógeno ($HCl$): Proviene de la descomposición del policloruro de vinilo ($PVC$) usado en
                    aislaciones de cables y cañerías. Con la humedad ambiente o pulmonar se convierte en ácido clorhídrico,
                    provocando quemaduras químicas en vías aéreas y corroyendo velozmente pistas de cobre y componentes en placas
                    electrónicas.
              \item Dióxido de azufre ($SO_2$): Desprendido en la combustión de cauchos vulcanizados (neumáticos, sellos de
                    gabinetes, juntas y correas de transmisión) y lana, debido al azufre empleado como agente de vulcanización.
                    Forma ácido sulfuroso y sulfúrico, generando intenso ardor ocular y bronquial.
              \item Acroleína ($C_3H_4O$): Se genera por la pirólisis de grasas y aceites lubricantes de maquinarias y
                    transformadores (deshidratación térmica del glicerol) y sobrecalentamiento de resinas y maderas. Posee un
                    efecto lacrimógeno extremo que induce blefaroespasmo reflejo (cierre involuntario de párpados), cegando a
                    los ocupantes e impidiendo localizar las vías de evacuación.
              \item Dióxido de nitrógeno y óxidos de nitrógeno ($NO_x$): Originados por la reacción térmica a alta
                    temperatura ($>1000\Celsius$) entre el nitrógeno atmosférico ($N_2$) del aire y el oxígeno, así como por la
                    pirólisis de lacas, resinas termoestables y materiales nitrogenados. Su inhalación destruye las membranas
                    alveolares y desencadena un edema agudo de pulmón retardado (acumulación de líquido letal entre 12 y 24
                    horas posteriores a la exposición).
          \end{itemize}
    \item Disminución Crítica del Nivel de Oxígeno: La combustión consume aceleradamente el comburente disponible. Mientras
          que el aire limpio posee un $21\percent$ de oxígeno en volumen, valores inferiores al $16\percent$ provocan
          taquicardia y mareos, y una concentración menor al $6\percent$ causa la detención respiratoria en pocos segundos.
    \item Temperatura de la Mezcla Gaseosa: La inhalación de gases a temperaturas superiores a $120\Celsius$ provoca
          quemaduras térmicas inmediatas en las vías aéreas superiores, causando la muerte por laringoespasmo u oclusión.
    \item Oscurecimiento Óptico y Desorientación: Las partículas sólidas de hollín en suspensión dispersan y absorben la
          luz, reduciendo drásticamente la visibilidad hacia las señales de salida y obligando a los ocupantes a frenar su
          marcha de escape.
\end{itemize}

\subsection{Clasificación y Diagnóstico Visual del Humo según su Coloración}

En la práctica pericial y en la evaluación operativa del incendio, el color del humo emitido constituye un indicador
físico del tipo de combustión y de los materiales involucrados:

\begin{itemize}
    \item Humos Blancos: Indican combustión completa o con abundancia de comburente de productos vegetales, forrajes o
          maderas ligeras, o bien una elevada presencia de vapor de agua procedente del secado del material o de la
          extinción hídrica.
    \item Humos Amarillos: Denotan la presencia de sustancias químicas industriales con azufre, ácido clorhídrico o
          ácido nítrico. Alertan sobre atmósferas extremadamente tóxicas, irritantes y letales por inhalación.
    \item Humos Grises: Característicos de la combustión incompleta de compuestos celulósicos, papeles, telas y fibras
          artificiales bajo condiciones de ventilación restringida o en fase latente con generación de brasas.
    \item Humos Negros Claros: Señalan la combustión de cauchos, carbón vegetal o mezclas combustibles intermedias.
    \item Humos Negros Oscuros y Densos: Revelan la combustión de hidrocarburos pesados, petróleo, solventes orgánicos,
          asfaltos, plásticos espumados y fibras acrílicas. Poseen una alta carga de partículas de carbono libre (hollín)
          que bloquea totalmente la visibilidad y evidencia una severa deficiencia local de comburente.
\end{itemize}

El entendimiento riguroso de estas limitaciones fundamenta la obligatoriedad de los sistemas de evacuación y
compartimentación exigidos por la legislación laboral, así como la adopción de tecnologías de detección ultraprecoz y
las alternativas de prevención y protección analizadas en los capítulos subsiguientes.

